\subsection{WAWPart: workload-aware weighted partitioning графовых данных}

Метод WAWPart \cite{wawpart} относится к классу workload-aware алгоритмов разбиения графовых данных и ориентирован на минимизацию стоимости выполнения запросов обхода графа в распределённых графовых базах данных. В отличие от ранее рассмотренных подходов (Schism и SWORD), где основное внимание уделяется либо совместному использованию вершин в запросах, либо сжатому представлению рабочей нагрузки, WAWPart фокусируется на явном моделировании структуры обходов графа и интенсивности переходов между вершинами.

Ключевая идея метода заключается в том, что стоимость распределённого выполнения запросов определяется не только тем, какие вершины используются совместно, но и тем, как именно они последовательно посещаются в процессе обхода графа. Поэтому основным объектом оптимизации становится не просто edge-cut, а weighted traversal-cut, основанный на частоте и структуре переходов между вершинами.

\subsubsection{Постановка задачи WAWPART}

Рассматривается граф данных

\begin{equation}
G = (V, E),
\end{equation}

где $V$ --- множество вершин, $E$ --- множество рёбер.

Также задано множество запросов

\begin{equation}
Q = \{q_1, q_2, \dots, q_m\},
\end{equation}

где каждый запрос представляет собой последовательность обхода графа (последовательность посещаемых вершин и переходов между ними).

Требуется найти разбиение множества вершин на $k$ партиций:

\begin{equation}
\Pi = \{P_1, P_2, \dots, P_k\},
\end{equation}

такое, что:

\begin{equation}
\bigcup_{i=1}^{k} P_i = V, \quad P_i \cap P_j = \emptyset, \quad i \neq j,
\end{equation}

и минимизируется стоимость выполнения рабочей нагрузки обходов графа:

\begin{equation}
C(Q) = \sum_{q \in Q} C(q) \rightarrow \min.
\end{equation}

Стоимость определяется числом и «стоимостью» межпартиционных переходов при выполнении обходов графа.

\textbf{Ключевая идея: переход от вершин к переходам}

В отличие от Schism и SWORD, где базовой единицей моделирования является совместное использование вершин в запросах, WAWPart использует переходы между вершинами как основную единицу анализа.

Пусть запрос $q \in Q$ задаёт последовательность:

\begin{equation}
q = (v_1, v_2, \dots, v_n).
\end{equation}

Тогда он индуцирует множество переходов:

\begin{equation}
T(q) = \{(v_i, v_{i+1}) \mid i = 1, \dots, n-1\}.
\end{equation}

Именно эти переходы являются основой построения workload-модели.

\subsubsection{Построение workload-графа как графа переходов}

WAWPart строит workload-граф

\begin{equation}
G_w = (V, E_w),
\end{equation}

где $E_w$ отражает интенсивность переходов между вершинами в рабочей нагрузке.

В отличие от классического edge-based представления, здесь каждое ребро $e=(u,v)$ интерпретируется как агрегированная статистика всех переходов $u \rightarrow v$ во всех запросах.

\textbf{Базовое вычисление веса ребра}

Наиболее простая форма веса определяется как количество наблюдаемых переходов:

\begin{equation}
w(u,v) = \sum_{q \in Q} \sum_{(x,y)\in T(q)} \mathbb{I}[x=u \land y=v],
\end{equation}

где $\mathbb{I}[\cdot]$ --- индикаторная функция.

Эта модель отражает абсолютную частоту использования перехода в рабочей нагрузке.

\textbf{Учёт частоты выполнения запросов}

Если различные запросы выполняются с разной частотой, вводится вес запроса $f(q)$:

\begin{equation}
w(u,v) = \sum_{q \in Q} f(q)\cdot \sum_{(x,y)\in T(q)} \mathbb{I}[x=u \land y=y].
\end{equation}

Таким образом, вклад каждого перехода масштабируется в зависимости от реальной интенсивности выполнения запросов.

Это позволяет учитывать не только структуру workload, но и его распределение во времени.

\textbf{Нормализация переходов и вероятностная модель}

Помимо абсолютных частот, WAWPart использует вероятностную интерпретацию переходов.

Для вершины $u$ вводится вероятность перехода к $v$:

\begin{equation}
P(u,v) = \frac{N(u,v)}{\sum_{x \in V} N(u,x)},
\end{equation}

где $N(u,v)$ --- число переходов из $u$ в $v$.

Тогда вес может быть определён как:

\begin{equation}
w(u,v) = N(u,v)\cdot P(u,v).
\end{equation}

Такая форма усиливает влияние устойчивых и структурно значимых переходов, снижая влияние случайных шумовых переходов.

\textbf{Позиционная значимость переходов в обходах графа}

Важной особенностью WAWPart является учёт позиции перехода внутри обхода графа.

Пусть переход $(v_i, v_{i+1})$ находится на позиции $i$ в запросе длины $n$. Тогда вводится весовой коэффициент:

\begin{equation}
\phi(i,n) = \frac{1}{1 + \alpha \cdot i},
\end{equation}

где $\alpha$ --- параметр затухания значимости.

Тогда итоговый вклад перехода определяется как:

\begin{equation}
w(u,v) = \sum_{q \in Q} f(q)\cdot \sum_{(v_i,v_{i+1})\in T(q)} \mathbb{I}[\cdot]\cdot \phi(i,n).
\end{equation}

Интуитивно это означает, что ранние и критические переходы в обходе могут иметь больший вклад в стоимость разбиения.

\textbf{Учёт стоимости межпартиционных переходов}

WAWPart вводит дополнительную модель стоимости пересечения границ партиций.

Если переход $(u,v)$ выполняется между разными партициями, он вносит штраф:

\begin{equation}
c(u,v) = \delta(u,v)\cdot w(u,v),
\end{equation}

где:

\begin{equation}
\delta(u,v)=
\begin{cases}
1, & \text{если } u \text{ и } v \text{ находятся в разных партициях},\\
0, & \text{иначе}.
\end{cases}
\end{equation}

Таким образом, итоговая цель разбиения сводится к минимизации:

\begin{equation}
\min \sum_{(u,v)\in E_w} w(u,v)\delta(u,v).
\end{equation}

\textbf{Пороговая фильтрация рёбер workload-графа}

Для уменьшения сложности обработки вводится фильтрация слабых переходов:

\begin{equation}
E_w = \{(u,v)\mid w(u,v) > \theta\}.
\end{equation}

Это позволяет:
\begin{itemize}
    \item исключить случайные переходы;
    \item уменьшить размер workload-графа;
    \item сфокусироваться на устойчивых структурах обхода графа.
\end{itemize}

\textbf{Интерпретация workload-графа}

Полученный workload-граф представляет собой не просто структуру совместного использования данных, а ориентированный граф интенсивности переходов при обходах графа.

Вершины соответствуют сущностям данных, а рёбра отражают вероятность и частоту переходов между ними в рамках реальных запросов.

Таким образом, workload-граф в WAWPart является моделью поведения системы, а не только её структуры.

\textbf{Итоговая формализация}

В результате построения workload-графа задача разбиения формулируется как оптимизация weighted traversal-cut:

\begin{equation}
\min \sum_{(u,v)\in E_w} w(u,v)\delta(u,v),
\end{equation}

при ограничении балансировки:

\begin{equation}
|P_i| \leq (1+\varepsilon)\frac{|V|}{k}.
\end{equation}

Таким образом, WAWPart сводит задачу распределения графовых данных к минимизации стоимости разрезания наиболее интенсивно используемых переходов в структуре обходов графа.

\subsubsection{Сравнение со Schism и SWORD}

По сравнению с Schism, метод WAWPart использует более специализированную модель рабочей нагрузки, ориентированную на запросы обхода графа. Если Schism моделирует совместное использование вершин в запросах, то WAWPart моделирует непосредственно переходы при обходе графа.

По сравнению с SWORD, WAWPart уделяет больше внимания локальности путей обхода и направлению переходов между вершинами.

Таким образом:
\begin{itemize}
    \item Schism ориентирован на совместное использование данных в запросах;
    \item SWORD ориентирован на масштабируемую аппроксимацию рабочей нагрузки;
    \item WAWPart ориентирован на оптимизацию локальности обходов графа.
\end{itemize}
